% =====================================================================
%  اتفاقية تزويد ملاعب وعرضها على منصة «المستديرة»
%  ---------------------------------------------------------------
%  الترجمة:  xelatex venue-agreement-ar.tex   (مرّتين، ليضبط الفهرس والمراجع)
%  يتطلّب:   XeLaTeX + polyglossia + bidi + خط عربي (Amiri مفضّلًا)
%
%  إن لم يكن Amiri مثبّتًا، بدّل اسم الخط في السطر المعلّم أدناه إلى خط
%  موجود على ويندوز، مثل:  Traditional Arabic  أو  Simplified Arabic  أو  Arial
%
%  ⚠ هذه صيغة عمل وليست استشارة قانونية. تُراجَع لدى محامٍ أردني قبل التوقيع.
% =====================================================================

\documentclass[12pt,a4paper]{article}

% ---------- الحزم (polyglossia آخرها دائمًا مع bidi) ----------
\usepackage{geometry}
\geometry{a4paper,top=2.4cm,bottom=2.4cm,left=2.2cm,right=2.2cm,footskip=1.2cm}

\usepackage{xcolor}
\usepackage{array}
\usepackage{tabularx}
\usepackage{booktabs}
\usepackage{enumitem}
\usepackage{setspace}
\usepackage{fancyhdr}
\usepackage{titlesec}
\usepackage[hidelinks]{hyperref}

\usepackage{fontspec}
\usepackage{polyglossia}
\setmainlanguage[numerals=mashriq]{arabic}
\setotherlanguage{english}

% <<<<<<<<<<  بدّل اسم الخط هنا إن لزم  >>>>>>>>>>
\newfontfamily\arabicfont[Script=Arabic,Scale=1.05]{Amiri}
\newfontfamily\arabicfontsf[Script=Arabic,Scale=1.05]{Amiri}
\newfontfamily\arabicfonttt[Script=Arabic,Scale=1.00]{Amiri}
\newfontfamily\englishfont[Scale=0.95]{Times New Roman}

% ---------- ألوان الهوية ----------
\definecolor{teal}{HTML}{0F4B53}
\definecolor{lime}{HTML}{8CC63E}
\definecolor{rule}{HTML}{C9D3D5}
\definecolor{ink}{HTML}{101820}

% ---------- إعدادات عامة ----------
\onehalfspacing
\setlength{\parindent}{0pt}
\setlength{\parskip}{5pt}
\color{ink}

% ---------- ترويسة وتذييل ----------
\pagestyle{fancy}
\fancyhf{}
\renewcommand{\headrulewidth}{0.4pt}
\renewcommand{\footrulewidth}{0.4pt}
\fancyhead[R]{\small\color{teal} اتفاقية تزويد ملاعب --- منصة «المستديرة»}
\fancyhead[L]{\small\color{teal} \today}
\fancyfoot[C]{\small صفحة \thepage\ من \pageref{LastPage}}
\fancyfoot[L]{\small توقيع الطرف الأول: \rule{2.2cm}{0.4pt}}
\fancyfoot[R]{\small توقيع الطرف الثاني: \rule{2.2cm}{0.4pt}}
\usepackage{lastpage}

% ---------- أوامر مساعدة ----------
\newcounter{bndc}
\newcommand{\bnd}[1]{%
  \vspace{10pt}\par\noindent
  {\color{teal}\bfseries\large المادة (\arabic{bndc}): #1}%
  \stepcounter{bndc}%
  \par\vspace{2pt}\noindent{\color{rule}\rule{\linewidth}{0.6pt}}\par\vspace{4pt}%
}
\setcounter{bndc}{1}

\newcommand{\blank}[1][3.2cm]{\rule[-0.4ex]{#1}{0.5pt}}
\newcommand{\mlh}[1]{{\bfseries #1}}

\newenvironment{bnditems}
  {\begin{enumerate}[label=\arabic{bndc}.\arabic*,leftmargin=1.6em,itemsep=3pt,topsep=3pt]}
  {\end{enumerate}}

\newenvironment{sublist}
  {\begin{enumerate}[label=(\alph*),leftmargin=1.6em,itemsep=2pt,topsep=2pt]}
  {\end{enumerate}}

\newcommand{\annex}[2]{%
  \newpage
  \begin{center}
    {\color{teal}\bfseries\Large الملحق (#1)}\\[2pt]
    {\bfseries\large #2}
  \end{center}
  \vspace{4pt}\noindent{\color{teal}\rule{\linewidth}{1pt}}\par\vspace{8pt}%
}

% =====================================================================
\begin{document}

% ---------------- صفحة الغلاف ----------------
\thispagestyle{empty}
\begin{center}
\vspace*{2.2cm}
{\color{teal}\bfseries\fontsize{30}{36}\selectfont المستديرة}\\[4pt]
{\color{lime}\rule{5cm}{2.5pt}}\\[26pt]
{\bfseries\fontsize{22}{28}\selectfont اتفاقية تزويد ملاعب}\\[8pt]
{\fontsize{16}{20}\selectfont وعرضها على منصة الحجز الإلكتروني}\\[36pt]

\begin{tabular}{@{}p{0.32\linewidth}@{}p{0.62\linewidth}@{}}
\mlh{رقم الاتفاقية} & : \blank[6cm]\\[8pt]
\mlh{تاريخ التوقيع}  & : \blank[6cm]\\[8pt]
\mlh{مكان التوقيع}   & : \blank[6cm]\\[8pt]
\mlh{اسم المزوّد}     & : \blank[6cm]\\[8pt]
\mlh{عدد الملاعب}    & : \blank[6cm]\\
\end{tabular}

\vfill
{\small\color{teal}
حُرّرت هذه الاتفاقية من نسختين أصليتين، تسلّم كلُّ طرفٍ نسخةً منها للعمل بموجبها.}\\[6pt]
{\scriptsize\color{rule}\rule{0.5\linewidth}{0.4pt}}
\end{center}
\newpage

% ---------------- الأطراف ----------------
\begin{center}
{\color{teal}\bfseries\Large الأطراف المتعاقدة}
\end{center}
\vspace{4pt}

حُرّرت هذه الاتفاقية في مدينة \blank[3cm] بتاريخ \blank[1.2cm]\,/\,\blank[1.2cm]\,/\,\blank[1.8cm]\,م، بين كلٍّ من:

\vspace{8pt}
\noindent{\color{teal}\bfseries أوّلًا: الطرف الأول (المنصة)}
\begin{tabular}{@{}p{0.30\linewidth}@{}p{0.66\linewidth}@{}}
الاسم / الكيان        & : \blank[8cm]\\[6pt]
الرقم الوطني للمنشأة  & : \blank[8cm]\\[6pt]
السجل التجاري / الرقم الضريبي & : \blank[8cm]\\[6pt]
العنوان               & : \blank[8cm]\\[6pt]
يمثّله                & : \blank[8cm]\\[6pt]
هاتف / بريد إلكتروني  & : \blank[8cm]\\
\end{tabular}

\vspace{4pt}
ويُشار إليه فيما بعد بـ\mlh{«الطرف الأول»} أو \mlh{«المنصة»}.

\vspace{12pt}
\noindent{\color{teal}\bfseries ثانيًا: الطرف الثاني (المزوّد)}
\begin{tabular}{@{}p{0.30\linewidth}@{}p{0.66\linewidth}@{}}
الاسم / الكيان         & : \blank[8cm]\\[6pt]
صفة الموقّع            & : مالك \; $\square$ \quad مفوّض بالتوقيع \; $\square$ \quad مستأجر مشغّل \; $\square$\\[6pt]
الرقم الوطني / السجل   & : \blank[8cm]\\[6pt]
اسم المنشأة الرياضية   & : \blank[8cm]\\[6pt]
العنوان (المدينة / المنطقة) & : \blank[8cm]\\[6pt]
هاتف                   & : \blank[8cm]\\[6pt]
بريد إلكتروني          & : \blank[8cm]\\
\end{tabular}

\vspace{4pt}
ويُشار إليه فيما بعد بـ\mlh{«الطرف الثاني»} أو \mlh{«المزوّد»}.

\vspace{6pt}
ويُشار إليهما مجتمعَين بـ\mlh{«الطرفين»}، وإلى كلٍّ منهما منفردًا بـ\mlh{«الطرف»}.

% ---------------- التمهيد ----------------
\vspace{14pt}
\noindent{\color{teal}\bfseries\large التمهيد}\par
\vspace{2pt}\noindent{\color{rule}\rule{\linewidth}{0.6pt}}\par\vspace{4pt}

\begin{itemize}[leftmargin=1.4em,itemsep=3pt,topsep=2pt]
  \item \mlh{لمّا كان} الطرف الأول يملك ويشغّل منصةً إلكترونيةً (تطبيق هاتف محمول وموقعًا إلكترونيًّا) تتيح للمستخدمين تصفّح الملاعب الرياضية وتقديم طلبات حجزٍ عليها؛
  \item \mlh{ولمّا كان} الطرف الثاني يملك أو يشغّل المنشأة الرياضية الموصوفة في \mlh{الملحق (أ)}، ويرغب في عرض ملاعبه على المنصة وتلقّي طلبات الحجز عبرها؛
  \item \mlh{وحيث} أقرّ الطرف الثاني بأهليته القانونية الكاملة للتعاقد، وبأنه يملك كامل الصلاحية للتصرّف بالملاعب موضوع هذه الاتفاقية وتأجيرها للغير؛
\end{itemize}

\noindent فقد اتفق الطرفان، وهما بكامل أهليتهما المعتبرة شرعًا وقانونًا، على ما يلي:

\newpage

% =====================  المواد  =====================

\bnd{التمهيد والملاحق}
\begin{bnditems}
  \item يُعدّ التمهيد أعلاه جزءًا لا يتجزّأ من هذه الاتفاقية ومكمّلًا لأحكامها ويُقرأ معها.
  \item تُعدّ الملاحق (أ) و(ب) و(ج) و(د) المرفقة جزءًا لا يتجزّأ من هذه الاتفاقية، ولها ذات القوة الإلزامية.
  \item في حال التعارض بين نصّ الاتفاقية ونصّ أيّ ملحق، يُعمل بنصّ \mlh{الاتفاقية} ما لم يكن الملحق قد عُدّل لاحقًا بموجب البند (\ref{art:amend}) وبتوقيع الطرفين.
\end{bnditems}

\bnd{التعريفات}
يكون للعبارات التالية، أينما وردت في هذه الاتفاقية، المعاني المبيّنة إزاءها ما لم يدلّ السياق على غير ذلك:

\begin{tabularx}{\linewidth}{@{}>{\bfseries\raggedright\arraybackslash}p{0.24\linewidth}@{\;:\;}X@{}}
\toprule
المنصة        & التطبيق والموقع الإلكتروني ولوحة التحكّم المملوكة للطرف الأول، بما فيها أنظمتها وقواعد بياناتها.\\
\midrule
الملعب        & أيّ ملعبٍ من الملاعب الواردة تفصيلًا في الملحق (أ)، المعروضة على المنصة.\\
\midrule
الخانة الزمنية & فترةٌ زمنيةٌ محدّدة البداية والنهاية يُتاح خلالها حجز الملعب، وفق الجدول المعتمد في الملحق (أ).\\
\midrule
المستخدم / اللاعب & أيّ شخصٍ يستخدم المنصة لتقديم طلب حجز.\\
\midrule
طلب الحجز     & الطلب الذي يقدّمه المستخدم عبر المنصة لحجز خانةٍ زمنيةٍ محدّدة، \mlh{ولا يُعدّ حجزًا مؤكّدًا} إلّا بقبول الطرف الثاني وفق البند (\ref{art:flow}).\\
\midrule
الحجز المؤكّد  & طلب الحجز بعد قبوله من الطرف الثاني عبر المنصة.\\
\midrule
الحجز المكتمل & الحجز المؤكّد الذي انقضى وقته دون إلغاءٍ من أيّ من الطرفين، سواء حضر المستخدم أم تخلّف عن الحضور.\\
\midrule
قيمة الحجز    & المبلغ الإجمالي المعروض على المنصة للخانة الزمنية، شاملًا كلّ الضرائب والرسوم.\\
\midrule
العمولة       & النسبة المستحقّة للطرف الأول من قيمة كل حجزٍ مكتمل، وفق الملحق (ب).\\
\midrule
فترة الردّ     & المدّة القصوى الملزِمة للطرف الثاني للردّ على طلب الحجز، وفق الملحق (ج).\\
\midrule
لوحة التحكّم   & الواجهة الإلكترونية التي يتيحها الطرف الأول للطرف الثاني لإدارة طلباته وبياناته.\\
\midrule
يوم عمل       & أيّ يومٍ عدا الجمعة والسبت والعطل الرسمية في المملكة الأردنية الهاشمية.\\
\bottomrule
\end{tabularx}

\bnd{موضوع الاتفاقية}
\begin{bnditems}
  \item يمنح الطرف الثاني الطرفَ الأول ترخيصًا غير حصريّ، قابلًا للإلغاء، بعرض ملاعبه وبياناتها وأسعارها وصورها على المنصة، وتلقّي طلبات الحجز عليها ونقلها إليه.
  \item يلتزم الطرف الأول بتوفير المنصة كوسيلةٍ تقنيةٍ للربط بين الطرف الثاني والمستخدمين، مقابل العمولة المبيّنة في الملحق (ب).
  \item لا تشمل هذه الاتفاقية أيّ التزامٍ على الطرف الأول بتحقيق عددٍ معيّنٍ من الحجوزات أو مستوى إيرادٍ معيّن، وأيّ تقديراتٍ يقدّمها هي تقديراتٌ غير ملزِمة.
\end{bnditems}

\bnd{طبيعة العلاقة بين الطرفين}\label{art:nature}
\begin{bnditems}
  \item العلاقة بين الطرفين هي علاقة \mlh{متعاقدَين مستقلَّين}. ولا تُنشئ هذه الاتفاقية بأيّ حالٍ من الأحوال علاقة شراكةٍ، أو مشروعٍ مشترك، أو وكالةٍ، أو عملٍ ووظيفة، أو تمثيلٍ قانوني بين الطرفين.
  \item الطرف الأول \mlh{وسيطٌ تقنيّ} ولا يُعدّ طرفًا في عقد استخدام الملعب المبرَم بين الطرف الثاني والمستخدم، ولا يملك الملعب ولا يشغّله ولا يشرف عليه ولا يقدّم الخدمة الرياضية فيه.
  \item يبقى الطرف الثاني هو المسؤول وحده، قِبَل المستخدم وقِبَل الغير، عن الملعب وسلامته وصلاحيته للاستعمال وعن الخدمة المقدَّمة فيه.
  \item لا يجوز لأيّ طرفٍ الالتزام باسم الطرف الآخر أو التعهّد نيابةً عنه أو الإيحاء بذلك.
\end{bnditems}

\bnd{التزامات الطرف الأول}
يلتزم الطرف الأول بما يلي:
\begin{bnditems}
  \item إتاحة المنصة وتشغيلها وصيانتها، وبذل العناية المعقولة لاستمرار عملها، وفق مستوى الخدمة المبيّن في \mlh{الملحق (ج)}.
  \item عرض ملاعب الطرف الثاني وبياناته وأسعاره كما زوّده بها، ونقل طلبات الحجز إليه دون تأخيرٍ غير مبرَّر.
  \item إتاحة لوحة تحكّمٍ للطرف الثاني تمكّنه من الاطلاع على طلباته وحجوزاته وقبولها أو رفضها.
  \item تزويد الطرف الثاني بكشف حسابٍ دوريّ بالعمولات المستحقّة وفق البند (\ref{art:comm}).
  \item منع الحجز المزدوج على الخانة الزمنية الواحدة على مستوى المنصة، بحيث لا تُعرض خانةٌ محجوزة أو معلّقة للحجز مرّةً أخرى.
  \item عدم تعديل أسعار الطرف الثاني أو أوقاته أو بياناته من جانبٍ واحد، عدا التصحيح الشكلي الواضح أو ما يقتضيه سحب بياناتٍ مخالفةٍ للقانون.
  \item المحافظة على سرّية بيانات الطرف الثاني وفق البند (\ref{art:conf}).
\end{bnditems}

\bnd{التزامات الطرف الثاني}
يلتزم الطرف الثاني بما يلي:
\begin{bnditems}
  \item \mlh{صحّة البيانات:} تزويد الطرف الأول ببيانات ملاعبه كاملةً وصحيحةً ومحدّثةً (المساحة، نوع الأرضية، المرافق، الإنارة، الأسعار، الخانات الزمنية)، وتحديثها خلال \mlh{٢٤ ساعة} من أيّ تغيير.
  \item \mlh{الترخيص:} أن يكون حائزًا على كافّة التراخيص والموافقات النظامية اللازمة لتشغيل الملعب وتأجيره، وأن يبقيها سارية طوال مدّة الاتفاقية، وأن يُعلم الطرف الأول فورًا بأيّ إيقافٍ أو سحبٍ لأيٍّ منها.
  \item \mlh{الردّ على الطلبات:} الردّ على كل طلب حجزٍ بالقبول أو الرفض عبر المنصة خلال فترة الردّ المحدّدة في الملحق (ج).
  \item \mlh{الوفاء بالحجز المؤكّد:} تخصيص الملعب والخانة الزمنية للمستخدم صاحب الحجز المؤكّد، وعدم تأجيرها للغير تحت أيّ ظرف.
  \item \mlh{عدم الحجز المزدوج:} تحديث توفّر الخانات على المنصة فور بيعها بأيّ قناةٍ أخرى، وتحمُّل المسؤولية الكاملة عن أيّ تعارضٍ ناشئ عن إخلاله بذلك.
  \item \mlh{وحدة السعر:} عدم مطالبة المستخدم القادم عبر المنصة بأيّ مبلغٍ يزيد على قيمة الحجز المعروضة، وعدم فرض أيّ رسومٍ إضافيةٍ غير معلَنةٍ عليها.
  \item \mlh{السلامة:} تهيئة الملعب ومرافقه بحالةٍ صالحةٍ وآمنةٍ للاستعمال، وتوفير الإنارة والمرافق المعلَن عنها.
  \item \mlh{عدم التمييز:} عدم رفض طلبات الحجز الواردة عبر المنصة لسببٍ تمييزيّ، أو بقصد توجيه المستخدم إلى قناةٍ أخرى.
  \item \mlh{الحساب:} المحافظة على سرّية بيانات دخوله إلى لوحة التحكّم، وتحمّل مسؤولية كل تصرّفٍ يقع من حسابه.
  \item \mlh{الإبلاغ:} إعلام الطرف الأول خلال يوم عملٍ واحدٍ بأيّ حادثٍ أو نزاعٍ مع مستخدمٍ قادمٍ عبر المنصة.
\end{bnditems}

\bnd{آلية الحجز والتأكيد والإلغاء}\label{art:flow}
\begin{bnditems}
  \item \mlh{طلبٌ لا التزام:} يُعدّ ما يقدّمه المستخدم عبر المنصة \mlh{طلب حجزٍ} لا حجزًا مؤكّدًا، ولا ينعقد عقد استخدام الملعب إلّا بقبول الطرف الثاني للطلب عبر المنصة.
  \item \mlh{حجز الخانة أثناء التعليق:} تُحجز الخانة الزمنية على المنصة ولا تُعرض للغير طوال فترة الردّ، ولذلك يلتزم الطرف الثاني بالردّ ضمن المهلة؛ إذ إنّ الطلب غير المُجاب عليه يمنع بيع الخانة ثمّ يضيع على الطرفين معًا.
  \item \mlh{انقضاء المهلة:} إذا انقضت فترة الردّ دون ردّ، يحقّ للطرف الأول إلغاء الطلب تلقائيًّا وإعلام المستخدم، دون أن يستحقّ الطرف الأول عمولةً عنه، ويُسجَّل ذلك في مؤشّر أداء الطرف الثاني.
  \item \mlh{إلغاء المستخدم:} يجوز للمستخدم إلغاء حجزه المؤكّد دون مقابلٍ قبل \mlh{ستّ (٦) ساعات} من بدء الخانة الزمنية، وبعد انقضاء هذه المهلة يتعذّر الإلغاء عبر المنصة ويتواصل المستخدم مع الطرف الثاني مباشرةً. ولا تستحقّ عمولةٌ عن الحجز الملغى.
  \item \mlh{إلغاء الطرف الثاني:} لا يجوز للطرف الثاني إلغاء حجزٍ مؤكّدٍ إلّا لسببٍ قاهرٍ أو فنّيٍّ جوهري (كعطبٍ يمنع اللعب)، وعليه إعلام الطرف الأول والمستخدم فورًا. وفي حال تجاوز إلغاءاته \blank[1.2cm]\% من حجوزاته المؤكّدة خلال شهرٍ واحد، جاز للطرف الأول تعليق عرض ملاعبه حتى معالجة السبب.
  \item \mlh{التخلّف عن الحضور:} إذا لم يحضر المستخدم دون إلغاء، يبقى الحجز \mlh{مكتملًا} وتستحقّ عليه العمولة، ويحقّ للطرف الثاني تحصيل قيمته وفق سياسته المعلَنة.
  \item \mlh{تعديل الموعد:} يجوز نقل الحجز إلى خانةٍ أخرى بموافقة الطرفين عبر المنصة، ويبقى الحجز المنقول محتفظًا بصفته لأغراض احتساب العمولة.
\end{bnditems}

\bnd{الأسعار وتحصيل قيمة الحجز}
\begin{bnditems}
  \item يحدّد الطرف الثاني منفردًا أسعار ملاعبه، ويلتزم بإدراجها على المنصة \mlh{شاملةً كلّ الضرائب والرسوم}.
  \item يجري تحصيل قيمة الحجز \mlh{نقدًا في الملعب} من المستخدم مباشرةً لصالح الطرف الثاني، ولا يتقاضى الطرف الأول أيّ مبلغٍ من المستخدم.
  \item الطرف الثاني هو الملتزم وحده بإصدار الفواتير الضريبية للمستخدمين وبأداء ما يترتّب عليه من ضرائب ورسوم عن قيمة الحجز.
  \item لأيّ من الطرفين اقتراح تعديل الأسعار، ويجري التعديل من لوحة التحكّم، ولا يسري إلّا على طلبات الحجز اللاحقة على إدخاله؛ أمّا الحجوزات المؤكّدة قبله فتبقى بسعرها الأصلي.
  \item إذا استُحدثت لاحقًا وسيلة دفعٍ إلكترونيةٍ عبر المنصة، عُدِّل هذا البند بملحقٍ موقَّعٍ يبيّن آلية التحصيل والتوريد ومدده.
\end{bnditems}

\bnd{العمولة واستحقاقها وتسويتها}\label{art:comm}
\begin{bnditems}
  \item يستحقّ الطرف الأول عمولةً بالنسبة المبيّنة في \mlh{الملحق (ب)} من قيمة كل \mlh{حجزٍ مكتمل} تمّ عبر المنصة.
  \item \mlh{لا تستحقّ العمولة} عن: طلب الحجز المرفوض، أو المنقضية مهلته، أو الحجز الملغى من المستخدم ضمن المهلة المسموحة، أو الملغى من الطرف الثاني، أو المتعذّر بقوّةٍ قاهرة.
  \item يصدر الطرف الأول \mlh{كشف حسابٍ} بالعمولات المستحقّة عن كل فترة تسوية (وفق الملحق (ب))، ويتيحه للطرف الثاني عبر لوحة التحكّم، ويتضمّن الكشف تفصيل الحجوزات وتواريخها وقيمها.
  \item \mlh{مدّة الاعتراض:} للطرف الثاني الاعتراض على الكشف كتابةً خلال \mlh{خمسة (٥) أيام} من إتاحته، مبيّنًا محلّ الاعتراض بالتحديد. وإذا انقضت المدّة دون اعتراضٍ عُدّ الكشف \mlh{مقبولًا ونهائيًّا} وحجّةً على الطرفين.
  \item \mlh{السداد:} يلتزم الطرف الثاني بسداد قيمة الكشف خلال \mlh{سبعة (٧) أيام} من تاريخ إتاحته، إلى الحساب البنكي المبيّن في الملحق (ب) أو بأيّ وسيلةٍ يتّفق عليها الطرفان كتابةً.
  \item \mlh{أثر التأخّر:} إذا تأخّر السداد أكثر من \mlh{سبعة (٧) أيام} على استحقاقه، جاز للطرف الأول \mlh{تعليق عرض ملاعب الطرف الثاني} على المنصة دون إشعارٍ إضافي وحتى السداد الكامل، دون أن يُخلّ ذلك بحقّه في المطالبة بالمبلغ وبأيّ تعويضٍ يستحقّه. ولا يُعدّ التعليق إنهاءً للاتفاقية.
  \item \mlh{المقاصّة:} للطرف الأول إجراء المقاصّة بين ما يستحقّه من عمولاتٍ وأيّ مبالغ يستحقّها الطرف الثاني قِبَله.
  \item \mlh{سجلّات المنصة:} تُعتمد سجلّات المنصة الإلكترونية بيّنةً في إثبات الحجوزات وأوقاتها وقيمها وحالاتها، وذلك عملًا بأحكام التشريعات الأردنية الناظمة للمعاملات الإلكترونية، ما لم يُقدَّم دليلٌ كتابيٌّ يخالفها.
  \item \mlh{الفترة الترويجية:} يجوز الاتفاق في الملحق (ب) على فترةٍ ترويجيةٍ تكون فيها العمولة صفرًا؛ وينتهي أثرها بانتهاء مدّتها تلقائيًّا دون حاجةٍ إلى إشعار.
\end{bnditems}

\bnd{عدم الالتفاف على المنصة}\label{art:bypass}
\begin{bnditems}
  \item يلتزم الطرف الثاني بعدم دفع المستخدمين الذين تعرّف عليهم \mlh{لأوّل مرّةٍ عبر المنصة} إلى إتمام حجوزاتهم اللاحقة خارجها بقصد تجنّب العمولة، سواء بالحسم المشروط أو بالطلب الصريح أو بأيّ وسيلةٍ مماثلة.
  \item إذا ثبت خلال \mlh{تسعين (٩٠) يومًا} من أوّل حجزٍ لمستخدمٍ عبر المنصة أنّ الطرف الثاني أتمّ معه حجزًا خارجها \mlh{بناءً على دعوةٍ أو تحريضٍ صادرٍ من الطرف الثاني}، استحقّت العمولة عن ذلك الحجز كما لو تمّ عبر المنصة.
  \item لا يُعدّ مخالفةً لهذا البند: حجز مستخدمٍ كان عميلًا للطرف الثاني قبل هذه الاتفاقية، أو حجزٌ بادر إليه المستخدم من تلقاء نفسه دون دعوةٍ من الطرف الثاني.
  \item عبء إثبات المخالفة على الطرف الأول.
\end{bnditems}

\bnd{عدم الحصرية}
\begin{bnditems}
  \item هذه الاتفاقية \mlh{غير حصرية}؛ فللطرف الثاني الاستمرار في تأجير ملاعبه بأيّ قناةٍ أخرى، وللطرف الأول التعاقد مع أيّ ملاعب أخرى بما فيها المنافسة.
  \item لا يُقيَّد الطرف الثاني في تسعير قنواته الأخرى، على أن يلتزم بأحكام البند (٨/٦) بشأن عدم تجاوز السعر المعروض على المنصة قِبَل مستخدميها.
\end{bnditems}

\bnd{التقييمات والمراجعات}
\begin{bnditems}
  \item تتيح المنصة للمستخدمين تقييم الملاعب بعد الحجز المكتمل، وتُنشر التقييمات الحقيقية دون تدخّلٍ من الطرفين.
  \item لا يجوز للطرف الأول حذف تقييمٍ صحيحٍ لصالح الطرف الثاني أو ضدّه، ويقتصر حقّه في الحذف على التقييمات المخالفة للقانون أو المسيئة أو غير المرتبطة بحجزٍ فعلي.
  \item لا يجوز للطرف الثاني إغراء المستخدمين بمقابلٍ لتقديم تقييماتٍ إيجابية، ولا التأثير عليهم لسحب تقييمٍ سلبي.
  \item للطرف الثاني حقّ الردّ العلني على أيّ تقييمٍ عبر لوحة التحكّم متى أُتيحت هذه الخاصية.
\end{bnditems}

\bnd{الملكية الفكرية والمحتوى}
\begin{bnditems}
  \item تبقى المنصة وبرمجياتها وتصاميمها وعلاماتها التجارية وقواعد بياناتها مملوكةً بالكامل للطرف الأول، ولا تمنح هذه الاتفاقية الطرف الثاني أيّ حقٍّ فيها عدا حقّ الاستعمال المحدود المبيّن هنا.
  \item يمنح الطرف الثاني الطرفَ الأول ترخيصًا غير حصريّ، ومجّانيًّا، وقابلًا للإلغاء، باستعمال \mlh{اسم منشأته وشعارها وصور ملاعبها} التي زوّده بها، لغرض عرضها على المنصة والموقع وموادّ التسويق المرتبطة بها.
  \item يقرّ الطرف الثاني بأنه يملك الحقوق على الصور والمواد التي زوّد بها الطرف الأول، أو أنه مرخّصٌ له باستعمالها ومنح الترخيص المذكور، ويضمن للطرف الأول أيّ مطالبةٍ تنشأ عن خلاف ذلك.
  \item ينتهي هذا الترخيص بانتهاء الاتفاقية، ويلتزم الطرف الأول بإزالة المواد خلال \mlh{خمسة عشر (١٥) يومًا}، عدا ما استُعمل في نُسخٍ أرشيفيةٍ أو سجلّاتٍ محاسبيةٍ يلزم الاحتفاظ بها.
\end{bnditems}

\bnd{حماية البيانات الشخصية}\label{art:data}
\begin{bnditems}
  \item بيانات المستخدمين المجموعة عبر المنصة مملوكةٌ للطرف الأول ويعالجها بصفته \mlh{المتحكّم بالبيانات}، وفق التشريعات الأردنية النافذة لحماية البيانات الشخصية.
  \item يتيح الطرف الأول للطرف الثاني من بيانات المستخدم \mlh{الحدّ الأدنى اللازم لتنفيذ الحجز} فقط (الاسم ورقم الهاتف وتفاصيل الحجز).
  \item يلتزم الطرف الثاني باستعمال هذه البيانات \mlh{لغرض تنفيذ الحجز حصرًا}، ويُحظر عليه صراحةً:
  \begin{sublist}
    \item استعمالها لأيّ غرضٍ تسويقيٍّ أو ترويجيّ؛
    \item بيعها أو مشاركتها أو نقلها إلى أيّ طرفٍ ثالث؛
    \item تجميعها في قوائم أو قواعد بياناتٍ مستقلّة.
  \end{sublist}
  \item يلتزم كلّ طرفٍ باتخاذ التدابير التقنية والتنظيمية المعقولة لحماية البيانات، وبإعلام الطرف الآخر خلال \mlh{اثنتين وسبعين (٧٢) ساعة} بأيّ اختراقٍ أو تسريبٍ يمسّ بيانات المستخدمين.
  \item تبقى أحكام هذا البند نافذةً بعد انتهاء الاتفاقية أو إنهائها.
\end{bnditems}

\bnd{السرّية}\label{art:conf}
\begin{bnditems}
  \item يلتزم كلّ طرفٍ بالمحافظة على سرّية المعلومات التجارية والمالية والتقنية التي يطّلع عليها بحكم هذه الاتفاقية، وبعدم إفشائها للغير دون موافقةٍ كتابيةٍ مسبقة.
  \item لا يشمل الالتزام: المعلومات العامة أصلًا، أو التي كانت لدى الطرف قبل التعاقد، أو التي يوجب القانون أو جهةٌ مختصّة إفشاءها.
  \item يسري هذا البند طوال مدّة الاتفاقية ولمدّة \mlh{سنتين (٢)} من تاريخ انتهائها.
\end{bnditems}

\bnd{الضمانات والإقرارات}
يقرّ الطرف الثاني ويضمن للطرف الأول ما يلي:
\begin{bnditems}
  \item أنه يملك الصفة والصلاحية الكاملة لإبرام هذه الاتفاقية ولتأجير الملاعب موضوعها.
  \item أنّ جميع البيانات التي زوّد بها الطرف الأول صحيحةٌ ودقيقةٌ وغير مضلِّلة.
  \item أنّ الملاعب مرخَّصةٌ ومطابقةٌ لاشتراطات السلامة العامة والدفاع المدني المعمول بها.
  \item أنه ليس طرفًا في أيّ نزاعٍ أو التزامٍ يحول دون تنفيذ هذه الاتفاقية.
\end{bnditems}

\bnd{المسؤولية وحدودها والتعويض}\label{art:liab}
\begin{bnditems}
  \item \mlh{مسؤولية الطرف الثاني:} يتحمّل الطرف الثاني وحده المسؤولية الكاملة عن كلّ ضررٍ أو إصابةٍ أو خسارةٍ تلحق بالمستخدمين أو الغير داخل الملعب أو مرافقه أو بسببها، وعن جودة الخدمة، وعن أيّ إخلالٍ منه بأحكام هذه الاتفاقية.
  \item \mlh{التعويض:} يلتزم الطرف الثاني بتعويض الطرف الأول وحمايته من أيّ مطالبةٍ أو دعوى أو غرامةٍ أو نفقةٍ (بما فيها أتعاب المحاماة المعقولة) تنشأ عن: تشغيل الملعب، أو إصاباتٍ وقعت فيه، أو عدم صحّة بياناته، أو إخلاله بالتراخيص، أو مخالفته للبند (\ref{art:data}).
  \item \mlh{حدود مسؤولية الطرف الأول:} لا يُسأل الطرف الأول عن: ما يقع داخل الملعب، ولا عن سلوك المستخدمين أو تخلّفهم عن الحضور أو امتناعهم عن الدفع، ولا عن الأرباح الفائتة أو الفرص الضائعة، ولا عن الأضرار غير المباشرة أو التبعية أيًّا كان سببها.
  \item \mlh{السقف:} لا تتجاوز المسؤولية الإجمالية للطرف الأول تجاه الطرف الثاني، مهما تعدّدت أسبابها، \mlh{مجموع العمولات التي قبضها فعليًّا من الطرف الثاني خلال الأشهر الثلاثة (٣) السابقة} على الواقعة المنشئة للمسؤولية.
  \item \mlh{انقطاع الخدمة:} لا يُسأل الطرف الأول عن انقطاعٍ مؤقّتٍ في المنصة ناشئٍ عن صيانةٍ أو عطلٍ فنّيٍّ أو انقطاعٍ في خدمات مزوّدي البنية التحتية، على أن يبذل العناية المعقولة لمعالجته وفق الملحق (ج).
  \item لا يحدّ هذا البند من المسؤولية عن الغشّ أو الخطأ الجسيم أو ما لا يجوز الاتفاق على تحديد المسؤولية عنه قانونًا.
\end{bnditems}

\bnd{السلامة والتأمين}
\begin{bnditems}
  \item يلتزم الطرف الثاني بالمحافظة على الملعب ومرافقه بحالةٍ آمنةٍ، وبإجراء الصيانة الدورية، وبإعلان تعليمات السلامة في مكانٍ ظاهر.
  \item يلتزم الطرف الثاني بإبرام والاحتفاظ بوثيقة \mlh{تأمين مسؤوليةٍ مدنيةٍ تجاه الغير} سارية المفعول طوال مدّة الاتفاقية، وبتزويد الطرف الأول بنسخةٍ عنها عند الطلب.\footnote{إن تعذّر ذلك عند التوقيع، يُمنح الطرف الثاني مهلةً مدّتها \blank[1.4cm] يومًا لاستيفائه، ويُثبت ذلك في الملحق (أ).}
  \item لا يُنشئ عدم مطالبة الطرف الأول بالوثيقة أيّ تنازلٍ عن هذا الالتزام.
\end{bnditems}

\bnd{مدّة الاتفاقية والتجديد}
\begin{bnditems}
  \item تسري هذه الاتفاقية لمدّة \mlh{اثني عشر (١٢) شهرًا} تبدأ من تاريخ توقيعها.
  \item تتجدّد تلقائيًّا لمددٍ مماثلةٍ ما لم يُشعر أحد الطرفين الآخر برغبته في عدم التجديد قبل \mlh{ثلاثين (٣٠) يومًا} على الأقلّ من انتهاء المدّة السارية.
\end{bnditems}

\bnd{الإنهاء}
\begin{bnditems}
  \item \mlh{الإنهاء بالإرادة المنفردة:} لأيّ من الطرفين إنهاء الاتفاقية دون إبداء أسبابٍ بإشعارٍ كتابيٍّ مسبقٍ مدّته \mlh{ثلاثون (٣٠) يومًا}.
  \item \mlh{الإنهاء الفوري:} لأيّ من الطرفين إنهاؤها فورًا وبإشعارٍ كتابيّ، دون إخلالٍ بحقّه في التعويض، في أيٍّ من الحالات التالية:
  \begin{sublist}
    \item إخلال الطرف الآخر إخلالًا جوهريًّا لم يُصلحه خلال \mlh{سبعة (٧) أيام} من إنذاره كتابةً؛
    \item تكرار الطرف الثاني عدم الوفاء بحجوزاتٍ مؤكّدة، أو مخالفته الجسيمة للبند (\ref{art:bypass}) أو البند (\ref{art:data})؛
    \item تأخّر الطرف الثاني في سداد العمولات أكثر من \mlh{ثلاثين (٣٠) يومًا}؛
    \item سحب ترخيص الملعب أو إغلاقه؛
    \item إشهار إفلاس أحد الطرفين أو تصفيته أو توقّفه عن أعماله.
  \end{sublist}
  \item \mlh{أثر الإنهاء:}
  \begin{sublist}
    \item يُرفع عرض ملاعب الطرف الثاني عن المنصة فورًا؛
    \item \mlh{تُنفَّذ الحجوزات المؤكّدة} القائمة قبل تاريخ الإنهاء، وتستحقّ عليها العمولة كاملةً؛
    \item تُصفّى الحسابات بين الطرفين خلال \mlh{خمسة عشر (١٥) يومًا} من تاريخ الإنهاء؛
    \item تبقى نافذةً بعد الإنهاء أحكام البنود (\ref{art:data}) و(\ref{art:conf}) و(\ref{art:liab}) و(\ref{art:law}).
  \end{sublist}
\end{bnditems}

\bnd{القوّة القاهرة}
\begin{bnditems}
  \item لا يُسأل أيّ طرفٍ عن التأخّر أو عدم التنفيذ الناشئ عن قوّةٍ قاهرةٍ خارجةٍ عن إرادته المعقولة، كالكوارث الطبيعية والأوبئة وقرارات السلطات وانقطاع الاتصالات أو الكهرباء العامّ.
  \item يلتزم الطرف المتأثّر بإعلام الآخر خلال \mlh{ثلاثة (٣) أيام} وببذل الجهد المعقول لتخفيف الأثر.
  \item إذا استمرّت القوّة القاهرة أكثر من \mlh{ستّين (٦٠) يومًا}، جاز لأيّ من الطرفين إنهاء الاتفاقية دون تعويض.
  \item لا تُعدّ صعوبة الطرف الثاني المالية قوّةً قاهرة.
\end{bnditems}

\bnd{تعديل الاتفاقية}\label{art:amend}
\begin{bnditems}
  \item لا يقع أيّ تعديلٍ على هذه الاتفاقية إلّا بمحرّرٍ خطّيٍّ موقَّعٍ من الطرفين.
  \item يُستثنى من ذلك تحديث بيانات الملاعب والأسعار والخانات الزمنية في الملحق (أ)، إذ تُعتمد التحديثات المدخَلة عبر لوحة التحكّم وتُعدّ جزءًا من الملحق دون حاجةٍ إلى توقيعٍ جديد.
  \item أيّ تعديلٍ على \mlh{نسبة العمولة} لا يسري إلّا بموافقةٍ خطّيةٍ من الطرفين وبإشعارٍ مسبقٍ مدّته \mlh{ثلاثون (٣٠) يومًا}، ولا يسري بأثرٍ رجعي.
\end{bnditems}

\bnd{التنازل والحوالة}
\begin{bnditems}
  \item لا يجوز لأيّ طرفٍ التنازل عن حقوقه أو التزاماته الناشئة عن هذه الاتفاقية للغير دون موافقةٍ خطّيةٍ مسبقةٍ من الطرف الآخر.
  \item يُستثنى تنازل الطرف الأول لصالح شركةٍ تابعةٍ له أو خَلَفٍ عامٍّ في حال الاندماج أو بيع النشاط، على أن يُشعر الطرف الثاني بذلك.
\end{bnditems}

\bnd{الإشعارات}
\begin{bnditems}
  \item تُوجَّه الإشعارات إلى العناوين المبيّنة في صدر هذه الاتفاقية، وتُعدّ صحيحةً إذا سُلّمت باليد، أو أُرسلت بالبريد المسجّل، أو بالبريد الإلكتروني، أو عبر لوحة التحكّم.
  \item يُعدّ الإشعار الإلكتروني واصلًا في \mlh{يوم العمل التالي} لإرساله ما لم يثبت خلاف ذلك.
  \item يلتزم كلّ طرفٍ بإعلام الآخر بأيّ تغييرٍ في عنوانه أو وسائل اتصاله خلال \mlh{سبعة (٧) أيام}، وإلّا عُدّ الإرسال إلى العنوان الأخير صحيحًا ومنتجًا لآثاره.
\end{bnditems}

\bnd{أحكام عامة}
\begin{bnditems}
  \item \mlh{استقلال البنود:} إذا قُضي ببطلان بندٍ أو تعذّر تنفيذه، بقيت سائر البنود صحيحةً ونافذة، ويُستبدل بالبند الباطل بندٌ صحيحٌ يقارب مقصده.
  \item \mlh{عدم التنازل الضمني:} لا يُعدّ تسامح أحد الطرفين أو سكوته عن مطالبةٍ تنازلًا عن حقّه فيها أو في غيرها لاحقًا.
  \item \mlh{الاتفاق الكامل:} تمثّل هذه الاتفاقية وملاحقها كامل ما اتّفق عليه الطرفان، وتحلّ محلّ كلّ اتفاقٍ أو تفاهمٍ أو وعدٍ سابقٍ، شفويًّا كان أم كتابيًّا.
  \item \mlh{العناوين:} وُضعت عناوين المواد للتيسير فقط ولا يُعتدّ بها في التفسير.
  \item \mlh{لغة الاتفاقية:} حُرّرت باللغة العربية، وهي اللغة المعتمدة في التفسير، وأيّ ترجمةٍ لها للاستئناس فقط.
  \item \mlh{النسخ:} حُرّرت من نسختين أصليتين بيد كلّ طرفٍ نسخة.
\end{bnditems}

\bnd{القانون الواجب التطبيق وتسوية النزاعات}\label{art:law}
\begin{bnditems}
  \item تخضع هذه الاتفاقية وتُفسَّر وفقًا لقوانين \mlh{المملكة الأردنية الهاشمية} النافذة.
  \item يسعى الطرفان لتسوية أيّ نزاعٍ ينشأ عن هذه الاتفاقية \mlh{ودّيًّا} خلال \mlh{ثلاثين (٣٠) يومًا} من تاريخ إشعار أحدهما الآخر بالنزاع كتابةً.
  \item إذا تعذّرت التسوية الودّية، يكون الاختصاص \mlh{لمحاكم عمّان النظامية}.\footnote{بديلٌ اختياري: «يُحال النزاع إلى التحكيم وفق قانون التحكيم الأردني، أمام محكّمٍ فردٍ يتّفق عليه الطرفان، ومكان التحكيم عمّان ولغته العربية، ويكون قرار المحكّم نهائيًّا وملزمًا». يُختار أحد المسارين ويُشطب الآخر عند التوقيع.}
\end{bnditems}

% ---------------- التواقيع ----------------
\vspace{18pt}
\noindent{\color{teal}\bfseries\large التواقيع}\par
\vspace{2pt}\noindent{\color{teal}\rule{\linewidth}{1pt}}\par\vspace{10pt}

\noindent
وإثباتًا لما تقدّم، وقّع الطرفان هذه الاتفاقية في التاريخ المبيّن في صدرها، بعد قراءتها وفهم مضمونها والموافقة على كافّة بنودها وملاحقها.

\vspace{20pt}
\noindent
\begin{tabularx}{\linewidth}{@{}X@{\hspace{1cm}}X@{}}
\mlh{الطرف الأول (المنصة)} & \mlh{الطرف الثاني (المزوّد)}\\[10pt]
الاسم: \blank[4.2cm]      & الاسم: \blank[4.2cm]\\[10pt]
الصفة: \blank[4.2cm]      & الصفة: \blank[4.2cm]\\[10pt]
التوقيع: \blank[4cm]      & التوقيع: \blank[4cm]\\[10pt]
الختم:                     & الختم:\\[34pt]
التاريخ: \blank[4cm]      & التاريخ: \blank[4cm]\\
\end{tabularx}

\vspace{22pt}
\noindent\mlh{الشاهد الأول:} الاسم \blank[4cm] \quad التوقيع \blank[3cm]\\[8pt]
\noindent\mlh{الشاهد الثاني:} الاسم \blank[4cm] \quad التوقيع \blank[3cm]

% =====================  الملاحق  =====================

\annex{أ}{بيانات المزوّد والملاعب والخانات الزمنية}

\noindent\mlh{١. بيانات المنشأة}
\vspace{4pt}

\begin{tabularx}{\linewidth}{@{}>{\bfseries}p{0.30\linewidth}@{\;:\;}X@{}}
\toprule
اسم المنشأة كما يُعرض        & \blank[8cm]\\
المدينة                      & \blank[8cm]\\
المنطقة                      & \blank[8cm]\\
العنوان التفصيلي             & \blank[8cm]\\
الإحداثيات (خط الطول/العرض)  & \blank[8cm]\\
هاتف الحجوزات                & \blank[8cm]\\
مسؤول التواصل                & \blank[8cm]\\
رقم الترخيص وجهة إصداره      & \blank[8cm]\\
وثيقة التأمين (رقمها/انتهاؤها)& \blank[8cm]\\
\bottomrule
\end{tabularx}

\vspace{12pt}
\noindent\mlh{٢. الملاعب المشمولة}
\vspace{4pt}

\begin{tabularx}{\linewidth}{@{}>{\centering\arraybackslash}p{0.05\linewidth}Xp{0.13\linewidth}p{0.13\linewidth}p{0.15\linewidth}@{}}
\toprule
\mlh{\#} & \mlh{اسم الملعب} & \mlh{الرياضة} & \mlh{الحجم} & \mlh{سعر الساعة (د.أ)}\\
\midrule
١ & \blank[4.5cm] & \blank[1.8cm] & \blank[1.8cm] & \blank[2.2cm]\\[6pt]
٢ & \blank[4.5cm] & \blank[1.8cm] & \blank[1.8cm] & \blank[2.2cm]\\[6pt]
٣ & \blank[4.5cm] & \blank[1.8cm] & \blank[1.8cm] & \blank[2.2cm]\\[6pt]
٤ & \blank[4.5cm] & \blank[1.8cm] & \blank[1.8cm] & \blank[2.2cm]\\[6pt]
٥ & \blank[4.5cm] & \blank[1.8cm] & \blank[1.8cm] & \blank[2.2cm]\\
\bottomrule
\end{tabularx}

\vspace{12pt}
\noindent\mlh{٣. الخانات الزمنية المعتمدة}
\vspace{4pt}

تُعتمد الخانات المدخَلة على لوحة التحكّم بتاريخ التوقيع، ويجوز تحديثها وفق البند (\ref{art:amend}/٢).

\vspace{6pt}
\begin{tabularx}{\linewidth}{@{}>{\bfseries}p{0.26\linewidth}@{\;:\;}X@{}}
\toprule
أيام العمل            & \blank[8cm]\\
ساعة الفتح / الإغلاق  & \blank[3cm] \quad إلى \quad \blank[3cm]\\
مدّة الخانة الواحدة    & \blank[8cm]\\
أوقات الذروة (إن وُجدت)& \blank[8cm]\\
\bottomrule
\end{tabularx}

\vspace{12pt}
\noindent\mlh{٤. المرافق المعلَنة}
\vspace{4pt}

\begin{tabularx}{\linewidth}{@{}XXXX@{}}
$\square$ إنارة & $\square$ مواقف & $\square$ غرف تبديل & $\square$ دورات مياه\\[5pt]
$\square$ مظلّة & $\square$ كافيتيريا & $\square$ مدرّجات & $\square$ كرات\\[5pt]
$\square$ مياه شرب & $\square$ إسعافات أوّلية & $\square$ واي فاي & $\square$ \blank[2.4cm]\\
\end{tabularx}

\annex{ب}{العمولة وشروط التسوية المالية}

\begin{tabularx}{\linewidth}{@{}>{\bfseries}p{0.34\linewidth}@{\;:\;}X@{}}
\toprule
نسبة العمولة                & \blank[2cm]\% من قيمة كل حجزٍ مكتمل\\
\midrule
الفترة الترويجية            & \blank[2cm] يومًا من تاريخ التوقيع، تكون العمولة خلالها \mlh{صفرًا}\\
\midrule
تاريخ بدء احتساب العمولة    & \blank[1.4cm]\,/\,\blank[1.4cm]\,/\,\blank[2cm]\\
\midrule
دورية التسوية               & $\square$ أسبوعية \quad $\square$ نصف شهرية \quad $\square$ شهرية\\
\midrule
موعد إصدار الكشف            & خلال \blank[1.6cm] أيام من انتهاء فترة التسوية\\
\midrule
مهلة السداد                 & \blank[1.6cm] أيام من تاريخ إتاحة الكشف\\
\midrule
مهلة الاعتراض               & \blank[1.6cm] أيام من تاريخ إتاحة الكشف\\
\midrule
الحدّ الأدنى للتحويل        & \blank[2cm] دينار (يُرحَّل ما دونه إلى الفترة التالية)\\
\bottomrule
\end{tabularx}

\vspace{14pt}
\noindent\mlh{بيانات السداد لصالح الطرف الأول}
\vspace{4pt}

\begin{tabularx}{\linewidth}{@{}>{\bfseries}p{0.30\linewidth}@{\;:\;}X@{}}
\toprule
اسم المستفيد     & \blank[8cm]\\
البنك / الفرع    & \blank[8cm]\\
رقم الحساب       & \blank[8cm]\\
الآيبان (IBAN)   & \blank[8cm]\\
محفظة إلكترونية  & \blank[8cm]\\
\bottomrule
\end{tabularx}

\vspace{14pt}
\noindent\mlh{ملاحظات}
\begin{itemize}[leftmargin=1.4em,itemsep=3pt]
  \item تُحتسب العمولة على \mlh{قيمة الحجز المعروضة على المنصة} شاملةً الضرائب.
  \item لا تُحتسب عمولةٌ عن الحالات المستثناة في البند (\ref{art:comm}/٢).
  \item أيّ ضريبةٍ مستحقّة على العمولة نفسها تُضاف إلى قيمتها ويؤدّيها الملتزم بها قانونًا.
\end{itemize}

\annex{ج}{مستوى الخدمة (SLA)}

\noindent\mlh{١. التزامات الطرف الأول}
\vspace{4pt}

\begin{tabularx}{\linewidth}{@{}X>{\centering\arraybackslash}p{0.26\linewidth}@{}}
\toprule
\mlh{المؤشّر} & \mlh{المستوى الملتزَم به}\\
\midrule
جاهزية المنصة شهريًّا (عدا الصيانة المجدولة) & لا تقلّ عن \blank[1.6cm]\%\\
\midrule
الإشعار المسبق بالصيانة المجدولة & \blank[1.6cm] ساعة\\
\midrule
الاستجابة لبلاغ عطلٍ حرج & خلال \blank[1.6cm] ساعة\\
\midrule
إتاحة كشف الحساب & وفق الملحق (ب)\\
\midrule
قنوات الدعم & هاتف / واتساب / بريد إلكتروني\\
\bottomrule
\end{tabularx}

\vspace{14pt}
\noindent\mlh{٢. التزامات الطرف الثاني}
\vspace{4pt}

\begin{tabularx}{\linewidth}{@{}X>{\centering\arraybackslash}p{0.26\linewidth}@{}}
\toprule
\mlh{المؤشّر} & \mlh{المستوى الملتزَم به}\\
\midrule
فترة الردّ على طلب الحجز (طلب لموعدٍ بعد ٢٤ ساعة فأكثر) & \blank[1.6cm] دقيقة\\
\midrule
فترة الردّ على طلبٍ لموعدٍ خلال ٢٤ ساعة & \blank[1.6cm] دقيقة\\
\midrule
نسبة قبول الطلبات شهريًّا & لا تقلّ عن \blank[1.6cm]\%\\
\midrule
نسبة إلغاء الحجوزات المؤكّدة من جانبه & لا تزيد على \blank[1.6cm]\%\\
\midrule
تحديث بيانات الأسعار والتوفّر & خلال \blank[1.6cm] ساعة من التغيير\\
\bottomrule
\end{tabularx}

\vspace{10pt}
\noindent\mlh{٣. أثر الإخلال بمستوى الخدمة}
\begin{itemize}[leftmargin=1.4em,itemsep=3pt]
  \item عند تجاوز الحدود أعلاه لشهرين متتاليين، يُوجّه الطرف الأول \mlh{إنذارًا كتابيًّا} مهلته \mlh{خمسة عشر (١٥) يومًا} للمعالجة.
  \item إن استمرّ الإخلال، جاز \mlh{تخفيض ترتيب ظهور} ملاعب الطرف الثاني على المنصة أو \mlh{تعليق عرضها} حتى المعالجة.
  \item لا يُعدّ ذلك إنهاءً للاتفاقية ولا تنازلًا عن حقّ الطرف الأول في الإنهاء وفق البند (١٩).
\end{itemize}

\annex{د}{نموذج كشف الحساب الدوري}

\noindent
\begin{tabularx}{\linewidth}{@{}>{\bfseries}p{0.24\linewidth}@{\;:\;}X@{}}
\toprule
المزوّد        & \blank[7cm]\\
فترة الكشف     & من \blank[2.4cm] إلى \blank[2.4cm]\\
تاريخ الإصدار  & \blank[7cm]\\
\bottomrule
\end{tabularx}

\vspace{12pt}
\begin{tabularx}{\linewidth}{@{}>{\centering\arraybackslash}p{0.06\linewidth}p{0.15\linewidth}Xp{0.14\linewidth}p{0.14\linewidth}@{}}
\toprule
\mlh{\#} & \mlh{التاريخ} & \mlh{الملعب / الخانة} & \mlh{قيمة الحجز} & \mlh{العمولة}\\
\midrule
١ & \blank[2cm] & \blank[3.4cm] & \blank[1.8cm] & \blank[1.8cm]\\[5pt]
٢ & \blank[2cm] & \blank[3.4cm] & \blank[1.8cm] & \blank[1.8cm]\\[5pt]
٣ & \blank[2cm] & \blank[3.4cm] & \blank[1.8cm] & \blank[1.8cm]\\[5pt]
٤ & \blank[2cm] & \blank[3.4cm] & \blank[1.8cm] & \blank[1.8cm]\\[5pt]
٥ & \blank[2cm] & \blank[3.4cm] & \blank[1.8cm] & \blank[1.8cm]\\
\midrule
\multicolumn{3}{@{}r@{}}{\mlh{عدد الحجوزات المكتملة}} & \multicolumn{2}{c}{\blank[3cm]}\\[4pt]
\multicolumn{3}{@{}r@{}}{\mlh{إجمالي قيمة الحجوزات}} & \multicolumn{2}{c}{\blank[3cm]}\\[4pt]
\multicolumn{3}{@{}r@{}}{\mlh{إجمالي العمولة المستحقّة}} & \multicolumn{2}{c}{\blank[3cm]}\\[4pt]
\multicolumn{3}{@{}r@{}}{\mlh{رصيد سابق / مقاصّة}} & \multicolumn{2}{c}{\blank[3cm]}\\[4pt]
\multicolumn{3}{@{}r@{}}{\mlh{الصافي الواجب السداد}} & \multicolumn{2}{c}{\blank[3cm]}\\
\bottomrule
\end{tabularx}

\vspace{14pt}
\noindent
\mlh{إقرار:} أقرّ باطّلاعي على هذا الكشف وموافقتي على ما ورد فيه، وأنّ عدم اعتراضي خلال المهلة المقرّرة في الملحق (ب) يجعله نهائيًّا وحجّةً عليّ وفق البند (\ref{art:comm}/٤).

\vspace{18pt}
\noindent
\begin{tabularx}{\linewidth}{@{}X@{\hspace{1cm}}X@{}}
\mlh{عن الطرف الأول} & \mlh{عن الطرف الثاني}\\[8pt]
الاسم: \blank[4.2cm] & الاسم: \blank[4.2cm]\\[10pt]
التوقيع: \blank[4cm] & التوقيع: \blank[4cm]\\[10pt]
التاريخ: \blank[4cm] & التاريخ: \blank[4cm]\\
\end{tabularx}

\end{document}
